\subsection{Identifikation eines Bauteils anhand des Schaltzeichens}
% Alt: Aufgabe 11
\Aufgabe{
  \label{Aufg11}
  Was für ein Bauelement wird mit folgendem Symbol gekennzeichnet?
  \begin{figure}[H]
    \centering
    \input{Tikz/src/FigSchaltzeichen.tex}\alttext{Die Abbildung zeigt ein Schaltzeichen für ein elektrisches Bauteil. Das Symbol besteht aus einem Kreis, aus dem drei Linien herausführen. An einer der Linien befindet sich ein kleiner Pfeil, der vom Kreis weg zeigt. Die drei Linien sind unterschiedlich angeordnet: Zwei Linien verlaufen fast gerade nach außen, die dritte Linie mit dem Pfeil ist leicht abgewinkelt.}
    \label{fig:FigSchaltzeichen}
  \end{figure}
}


\Loesung{
  Beim dargestellten Bauelement handelt es sich um einen npn-Transistor.
  Siehe: Abschnitt \ref{fig:UebersichtBipolartransistoren}
}